\begin{recetaConImagen}{Fideos con Carne}{4 personas}{35 min activos · 1 h 30 min total}{receta:fideos_carne}
    {imagenes/fideos_carne.png}{Esta fue la primera vez que hice los fideos en el piso nuevo, ¡salieron de 10!}

    \begin{minipage}[t]{0.35\textwidth}
        \section*{Ingredientes}
        \begin{itemize}[leftmargin=*]
            \item 500 g de carne (pollo y costilla de cerdo)
            \item 1 cebolla
            \item 2 tomates maduros
            \item 2 pimientos (rojo y verde)
            \item 200 ml de vino blanco
            \item 1 hoja de laurel
            \item colorante o azafrán
            \item sal
            \item 1,5 l de caldo de pollo, aprox.
            \item 300 g de fideos gruesos
        \end{itemize}
    \end{minipage}
    \hfill
    \begin{minipage}[t]{0.6\textwidth}
        \section*{Preparación}
        \begin{enumerate}
            \item \textbf{Sellar la carne:} En una olla con aceite, dorar los trozos de pollo y costilla. 
            \item \textbf{Sofrito:} Añadir por orden: 1º la cebolla, 2º los pimientos, 3º el tomate. 
            \item \textbf{Reducción:} Cuando no haya caldo del sofrito, verter un vaso de vino blanco y esperar a que reduzca. 
            \item \textbf{Cocción:} Añadir caldo de pollo y/o agua hasta cubrir, el laurel, el colorante y la sal cuando hierva. Cocer a fuego suave durante 45-60 minutos, hasta que la carne esté tierna.
            \item \textbf{Finalizar:} Añadir los fideos y cocer otros 15-20 minutos, hasta que estén hechos.
        \end{enumerate}
    \end{minipage}

    \vspace{1em}
    \begin{tcolorbox}[colback=orange!5, colframe=orange!50, title=Sugerencia, size=small]
        El vaso de vino no es obligatorio, pero mejora el sabor. 
        
        \vspace{0.5em} % Esto crea el salto de línea entre frases
        Un poco de queso en polvo por encima también le va muy bien. 
    \end{tcolorbox}
\end{recetaConImagen}
